% =====================================================================
% System Implementation (Weeks 5-8) -- paste into main.tex where the
% Month 2 / implementation section belongs. Uses the same \cite keys as
% the Review 1 paper (adjust to your actual BibTeX keys if they differ:
% gorkemli2015 -> IEEE 7130421, shahriar2024 -> arXiv:2403.15975,
% deo2024 -> PeerJ CS 2024, serag2025 -> Springer JNSM 2025).
%
% Preamble requirements (add once to main.tex if not already present):
%   \usepackage{algorithm}
%   \usepackage{algpseudocode}   % or \usepackage{algorithmic} for older style
% =====================================================================

\section{System Implementation}
\label{sec:implementation}

Having established the baseline degradation of real-time traffic under
contention in Section~\ref{sec:baseline}\footnote{Reference to the Review 1
baseline measurement section; rename to match your actual section label.},
this section describes the implementation of the four-layer pipeline that
addresses it: an SDN priority engine, a behavioral traffic classifier, the
integration bridge connecting the two, and a live monitoring dashboard.

\subsection{SDN-Based Priority Engine}
\label{subsec:priority-engine}

The priority engine is implemented as an os-ken\footnote{os-ken is the
actively maintained fork of the now-unmaintained Ryu SDN framework.}
controller application operating over OpenFlow~1.3 on an Open vSwitch
instance within a Mininet-emulated topology. Rather than a strict
precedence queue -- shown by Gorkemli et al.~\cite{gorkemli2015} to starve
lower-priority flows entirely under sustained contention -- we implement a
three-tier, fairness-aware scheme with a guaranteed minimum-bandwidth floor
per tier, following the bandwidth-splitting formulation of
Shahriar et al.~\cite{shahriar2024}.

Each tier is bound to a Linux-HTB queue on the switch's egress port:

\begin{itemize}
  \item \textbf{Tier 0 (Real-time):} minimum 6~Mbps, maximum 10~Mbps, highest scheduling priority.
  \item \textbf{Tier 1 (Best-effort):} minimum 2~Mbps, maximum 8~Mbps.
  \item \textbf{Tier 2 (Bulk):} minimum 1~Mbps (guaranteed floor), maximum 10~Mbps.
\end{itemize}

Algorithm~\ref{alg:bandwidth-split} summarizes the bandwidth-allocation
logic implemented by \texttt{priority\_controller.py}, adapted from
Algorithm~2 of Shahriar et al.~\cite{shahriar2024}.

\begin{algorithm}[t]
\caption{Tiered Bandwidth Allocation (adapted from \cite{shahriar2024})}
\label{alg:bandwidth-split}
\begin{algorithmic}[1]
\REQUIRE Flow $f$, predicted tier $t \in \{\text{realtime}, \text{besteffort}, \text{bulk}\}$, link capacity $C$
\STATE $q \gets \textsc{QueueForTier}(t)$
\IF{priority engine is enabled}
    \STATE install OpenFlow rule: match $f$, action \texttt{SetQueue}($q$), output normal
\ELSE
    \STATE install OpenFlow rule: match $f$, action output normal \COMMENT{no differentiation, reproduces baseline}
\ENDIF
\STATE guarantee $\text{min\_rate}(q) > 0$ for every $q$ \COMMENT{starvation-safe floor}
\end{algorithmic}
\end{algorithm}

A REST interface exposed by the controller (\texttt{/classify},
\texttt{/toggle}, \texttt{/status}) decouples flow-priority decisions from
rule installation, so that the classifier (Section~\ref{subsec:classifier})
and dashboard (Section~\ref{subsec:dashboard}) can drive the engine without
depending on OpenFlow internals directly.

\subsection{Behavioral Traffic Classifier}
\label{subsec:classifier}

Classification is performed without manual tagging or static port/IP
rules, addressing the limitation of static IP/port-based prioritization
identified by Deo et al.~\cite{deo2024}. For each flow, a rolling window of
20 packets is used to compute five behavioral features: mean and standard
deviation of packet size, mean and standard deviation of inter-arrival
time, and a burstiness ratio (coefficient of variation of inter-arrival
time). These features follow the classification pipeline of
Serag et al.~\cite{serag2025}, who show that packet-size and timing
statistics alone are sufficient to distinguish interactive from bulk
traffic without deep packet inspection.

A shallow decision tree (maximum depth 4) is trained on labeled flow
samples and used for live prediction; an interpretable rule-based fallback
is retained for cases where a trained model is unavailable, ensuring the
pipeline degrades gracefully rather than failing silently.

\subsection{Integration Pipeline}
\label{subsec:integration}

The integration layer connects the classifier's live output to the
controller's REST interface: as each flow accumulates a full feature
window, its predicted tier is pushed to \texttt{/classify}, which installs
the corresponding OpenFlow rule with no human intervention. This closes
the loop between detection (Section~\ref{subsec:classifier}) and
enforcement (Section~\ref{subsec:priority-engine}) that, per our gap
synthesis in the literature review, no single prior system combines
end-to-end.

\subsection{Live Monitoring Dashboard}
\label{subsec:dashboard}

A React/Node.js dashboard subscribes over WebSocket to a rolling metrics
feed (per-tier jitter, derived from the same inter-arrival statistics used
for classification) and exposes a single control: an ON/OFF toggle for the
priority engine. Toggling the engine calls the controller's
\texttt{/toggle} endpoint, which re-installs flows either with per-tier
queue differentiation (engine ON) or with uniform best-path forwarding
(engine OFF) -- allowing the baseline-vs-QoS comparison from
Section~\ref{sec:baseline} to be reproduced live, on demand, for
demonstration purposes.
